\subsection{The original canonical Cartier module map under geometric admissibility}
\begin{lemma}[Every actual principal Cartier equation is regular]
\label{lem:cartier-principal-equation-regular}
\lean{AlgebraicGeometry.Scheme.IdealSheafData.IsInvertible.generator_mem_nonZeroDivisors}
\leanok
Let $X$ be ANY original scheme, $J$ an actual effective Cartier ideal,
$U$ an actual affine open and $g\in\Gamma(U)$ with $J(U)=(g)$.
Then $g$ is a nonzerodivisor in the ORIGINAL section ring $\Gamma(U)$.
No reducedness, integrality, nonemptiness, Noetherian or finiteness
hypothesis is imposed. The empty affine and zero ring are included.
\end{lemma}
\begin{proof}
It suffices to show that $gt=0$ implies $t=0$. At each original $x\in U$,
choose an actual affine Cartier neighborhood $W$ with original regular
equation $h$. Choose an affine basic refinement $V\subseteq U\cap W$
containing $x$. The restriction of $h$ from $W$ to $V$ is regular:
the genuine map between affine section rings induced by the original
open inclusion is flat, and regularity is preserved by flat scalar
extension. The original ideal restriction formulas give
$(h|_V)=J(V)=(g|_V)$. In particular $h|_V=a(g|_V)$ for an actual
$a\in\Gamma(V)$. Restricting $gt=0$ therefore gives
$(h|_V)(t|_V)=a(g|_V)(t|_V)=0$, and regularity gives $t|_V=0$.
The original structure sheaf's section-extensionality over these actual
neighborhoods gives $t=0$. The right-cancellation statement follows
from commutativity. When $U$ is empty, section-extensionality is vacuous
and gives the correct zero-ring regularity statement.
\end{proof}

\begin{theorem}[Geometric Cartier admissibility makes the SAME canonical map an isomorphism]
\label{thm:admissible-cartier-canonical-module-pullback}
\lean{AlgebraicGeometry.Scheme.IdealSheafData.IsInvertible.toModulesComapPullbackMap_isIso_of_isInvertible_comap}
\leanok
\uses{lem:cartier-principal-equation-regular,lem:nonflat-cartier-principal-tensor,lem:affine-cartier-module-pullback-tensor,lem:cartier-ideal-locally-free-module,lem:actual-ideal-module-canonical-pullback-comp-iso,lem:actual-ideal-open-canonical-pullback-iso,scheme:genuine-ideal-pullback-stalk}
Let $f:X\to Y$ be ANY original scheme morphism. Let $I$ be an actual
effective Cartier ideal on $Y$ and assume the ACTUAL geometric ideal
comap $I\mathcal O_X$ is effective Cartier. Then the EXISTING canonical
comparison
\[f^*(I.\mathcal O_Y)\longrightarrow (I\mathcal O_X).\mathcal O_X\]
(the original \texttt{toModulesComapPullbackMap}) is an actual whole
module-sheaf isomorphism. No flatness, dominance, reducedness,
Noetherian, finite type, separation or quasi-compactness assumption
is imposed on either original scheme or on $f$.
The input is only geometric effective-Cartier admissibility; no map
injectivity, isomorphism, tensor identity or image certificate is assumed.
\end{theorem}
\begin{proof}
Fix an actual $x\in X$. Choose a true affine Cartier equation $g$ for $I$
on $U\ni f(x)$, and an original affine $V\ni x$ inside $f^{-1}U$.
Write $r:V\to U$ for the original restriction. The original Cartier ideal
$I_U$ has its original regular principal equation on the whole affine $U$.
The actual ideal comap $I_U\mathcal O_V$ equals the actual restriction
of $I\mathcal O_X$ to $V$, by the original comap-composition law and
the original coefficient square
$r\gg(U\hookrightarrow Y)=(V\hookrightarrow X)\gg f$.
Thus it is Cartier, since restriction along the original open immersion
is flat. Its ideal of whole sections is generated by the ACTUAL image
$g_V$ under the original coefficient map, by the genuine ideal-comap formula.
The preceding lemma proves $g_V$ regular, even on a nonreduced $V$.

Use the original affine pullback unit/tensor comparison, the original
ideal-module sections equivalences and the proved nonflat regular-principal
canonical tensor equivalence. On every actual pure tensor the resulting
map has included value $s\,r^\sharp(a)$, by the EXISTING canonical
\texttt{toModulesComapPullbackMap\_affine\_tmul} formula. The nonflat tensor
map has exactly the same value. Tensor induction and the literal ideal
inclusion give equality on the WHOLE original affine pullback module.
Thus the SAME actual affine comparison is bijective on whole sections.
Its source is quasicoherent because the original Cartier module is
locally free and genuine affine pullback preserves quasicoherence; its
target is locally free because the actual comap is Cartier. The EXISTING
quasicoherent affine isomorphism criterion promotes the actual whole-section
bijection to an isomorphism of the actual sheaf comparison for $r$.

Finally reuse the EXISTING canonical composition identity and cancellation
with open-immersion comparison isomorphisms to identify the actual
restriction of the original comparison for $f$ on $V$ with that genuine
affine isomorphism. The actual source opens $V$ constructed for every
$x$ cover $X$. The EXISTING module-sheaf restriction isomorphism criterion
proves the original canonical global comparison an isomorphism.
No replacement module, new geometric ideal, reduction, arbitrary chosen
isomorphism or compatibility hypothesis has entered the conclusion.
\end{proof}

\begin{corollary}[Dominant integral-source canonical Cartier comparison]
\label{cor:dominant-integral-cartier-canonical-module-pullback}
\lean{AlgebraicGeometry.Scheme.IdealSheafData.IsInvertible.toModulesComapPullbackMap_isIso_of_dominant_integral}
\leanok
\uses{thm:admissible-cartier-canonical-module-pullback,lem:nonflat-dominant-cartier-comap}
For ANY actual dominant $f:X\to Y$ with original $X$ integral and
original $Y$ reduced, and every actual effective Cartier ideal $I$ on $Y$,
the original canonical ideal-module comparison for $f$ is an isomorphism.
No quasi-compactness or flatness hypothesis on $f$ is needed.
\end{corollary}
\begin{proof}
The already proved original dominant integral-source Cartier-comap theorem
derives actual geometric admissibility. Apply the preceding theorem.
\end{proof}

\begin{corollary}[Original normalization canonical Cartier comparison]
\label{cor:normalization-cartier-canonical-module-pullback}
\lean{EnvelopeProperty.IntegralSchemeNormalization.cartierModuleComparison_isIso}
\leanok
\uses{thm:admissible-cartier-canonical-module-pullback,lem:nonflat-normalization-cartier-comap}
For ANY integral original scheme $Y$, the original generic-function-field
normalization projection $n:Y^\nu\to Y$ and ANY original effective Cartier
ideal $I$ induce an isomorphism of the SAME canonical ideal-module comparison.
No finite normalization, Noetherian, excellent, field, finite type,
quasi-compact, separated or flat hypothesis is imposed on original $Y$.
\end{corollary}
\begin{proof}
The already proved theorem for this literal normalization derives the
Cartier property of the actual comap. Apply the general canonical comparison.
This is a genuine normalization/line-bundle comparison prerequisite; the
whole normalized polarized model, its group action, relative-nef
approximation and LocalPositivity theorem remain separate obligations.
\end{proof}
